\documentclass[12pt]{article}

% === PACKAGES ===
\usepackage{amsmath, amssymb, amsthm} % For advanced math typesetting
\usepackage{graphicx}                 % For importing python graphs
\usepackage{geometry}                 % For page margins
\usepackage{hyperref}                 % For clickable links and cross-referencing
\usepackage{float}                    % For strict image placement [H]
\geometry{a4paper, margin=1in}

% === TITLE BLOCK ===
\title{Survival of the Fittest: A Mathematical Study of Multi-Strain SIR Dynamics under Differential Vaccine Efficacy}
\author{Muhammad Saad, Meesum Abbas}
\date{\today}

\begin{document}

\maketitle

% ==========================================
% 1. ABSTRACT
% ==========================================
% ==========================================
% 1. ABSTRACT
% ==========================================
\begin{abstract}
This project investigates the competitive dynamics of two distinct viral strains sharing a single host population. We extend the standard epidemiological framework to an SVIR (Susceptible-Vaccinated-Infected-Recovered) model to analyze how differential vaccine efficacy and varying transmission rates dictate which strain ultimately dominates. To establish the mathematical foundations of the model, we perform an in-depth analysis of both the local and global stability of the system, determining the exact conditions under which the viruses either die out naturally or persist in the population. Computationally, we simulate the system over a six-month period to observe the real-time phase transitions and displacement of the variants. Our analytical and numerical results confirm the Competitive Exclusion Principle, demonstrating that a highly infectious and vaccine-evasive strain will inevitably outcompete a lesser strain, driving it to extinction even in the presence of a mass vaccination campaign.
\end{abstract}

% ==========================================
% 2. INTRODUCTION
% ==========================================
% ==========================================
% 2. INTRODUCTION AND MATHEMATICAL MODEL
% ==========================================
\section{Introduction and Mathematical Model}

In the natural world, viruses are constantly mutating. When two distinct strains of the same virus circulate simultaneously, they do not just infect people independently; they actively compete for the same pool of vulnerable hosts. A prime real-world example of this occurred during the COVID-19 pandemic, when the highly infectious Omicron variant rapidly displaced the older Delta variant. 

This dynamic becomes even more complex when a population is vaccinated. If a vaccine is highly effective against the older strain but ``leaky" against the newer, mutated strain, the vaccine actually acts as an evolutionary filter, giving the new strain a massive competitive advantage. 

To explore this mathematically, we extend the classic SIR (Susceptible-Infected-Recovered) epidemiological model to an \textbf{SVIR} model featuring two competing viral strains. We divide the total population into five distinct compartments:
\begin{itemize}
    \item \textbf{$S$ (Susceptible):} Individuals with no immunity who can be infected by either strain.
    \item \textbf{$V$ (Vaccinated):} Individuals who have received a vaccine, granting them partial protection.
    \item \textbf{$I_1$ (Infected by Strain 1):} Individuals currently infected by the baseline strain (e.g., Delta).
    \item \textbf{$I_2$ (Infected by Strain 2):} Individuals currently infected by the mutated, evasive strain (e.g., Omicron).
    \item \textbf{$R$ (Recovered):} Individuals who have survived infection and gained natural immunity. 
\end{itemize}

We model the flow of people between these compartments over time ($t$) using the following system of non-linear ordinary differential equations:

\begin{align}
\frac{dS}{dt} &= \Lambda - (\beta_1 I_1 + \beta_2 I_2)S - (\nu + \mu)S \label{eq:S} \\
\frac{dV}{dt} &= \nu S - (\sigma_1 \beta_1 I_1 + \sigma_2 \beta_2 I_2)V - \mu V \label{eq:V} \\
\frac{dI_1}{dt} &= (\beta_1 S + \sigma_1 \beta_1 V)I_1 - (\gamma_1 + \mu)I_1 \label{eq:I1} \\
\frac{dI_2}{dt} &= (\beta_2 S + \sigma_2 \beta_2 V)I_2 - (\gamma_2 + \mu)I_2 \label{eq:I2} \\
\frac{dR}{dt} &= \gamma_1 I_1 + \gamma_2 I_2 - \mu R \label{eq:R}
\end{align}

To understand how the disease spreads, we define the following parameters:
\begin{itemize}
    \item \textbf{$\Lambda$ (Birth/Recruitment Rate):} The rate at which new, susceptible individuals enter the population.
    \item \textbf{$\mu$ (Natural Mortality Rate):} The rate at which people die from natural causes entirely unrelated to the virus.
    \item \textbf{$\nu$ (Vaccination Rate):} The speed at which susceptible people are being vaccinated.
    \item \textbf{$\beta_1, \beta_2$ (Transmission Rates):} The inherent infectivity of Strain 1 and Strain 2. This represents how easily the virus jumps from an infected person to a susceptible person.
    \item \textbf{$\gamma_1, \gamma_2$ (Recovery Rates):} The rate at which infected individuals recover and move into the $R$ compartment.
    \item \textbf{$\sigma_1, \sigma_2$ (Vaccine Leakage):} This is a critical parameter representing how well the virus evades the vaccine. It is defined as $\sigma = 1 - \epsilon$, where $\epsilon$ is the vaccine's efficacy. For example, if a vaccine is 90\% effective ($\epsilon = 0.90$), the leakage is 10\% ($\sigma = 0.10$). A leakage of $1.0$ means the vaccine provides zero protection.
\end{itemize}

\textbf{Core Assumptions:} For this model, we assume a closed, stable environment. The total population is normalized to a fraction of $N=1$, meaning $S, V, I_1, I_2,$ and $R$ represent percentages of the total population rather than raw headcounts. To maintain this constant population size, we assume that the birth rate perfectly balances the natural death rate ($\Lambda = \mu$). Consequently, the sum of all compartments is always $S + V + I_1 + I_2 + R = 1$ at any given time.

% ==========================================
% 3. LITERATURE REVIEW
% ==========================================
\section{Literature Review}

The mathematical modeling of infectious diseases has evolved significantly since the early 20th century, growing from simple compartmental models to highly complex systems that account for human behavior, medical interventions, and viral evolution. To ground our multi-strain SVIR model, we build upon several foundational texts and landmark papers in mathematical epidemiology.

The bedrock of our approach is derived from the classic textbook \textit{Mathematical Models in Population Biology and Epidemiology} by Fred Brauer and Carlos Castillo-Chavez. Their work exhaustively details the transition from the standard SIR model to models incorporating vital dynamics (births and deaths) and vaccination compartments. They demonstrated that while basic vaccination lowers the susceptible pool, the introduction of waning immunity or vaccine ``leakage" severely complicates the mathematical threshold for achieving herd immunity \hyperref[brauer2012]{(Brauer and Castillo-Chavez 45)}. 

The ultimate analytical goal of these complex models is to calculate the \textbf{Basic Reproduction Number ($R_0$)}. Conceptually, $R_0$ represents the average number of new infections caused by a single infected individual entering a completely susceptible population. Finding this number is the holy grail of epidemiology because it acts as a strict mathematical threshold: if $R_0 > 1$, the disease will explode into an epidemic; if $R_0 < 1$, the disease will naturally die out. 

Historically, calculating $R_0$ for models with multiple infected compartments (like our two-strain model) required solving agonizingly complex polynomial equations. To solve this, researchers Pauline van den Driessche and James Watmough published a landmark paper introducing the Next-Generation Matrix (NGM) method\footnote{\textbf{Next-Generation Matrix (NGM):} A mathematical shortcut used in systems with multiple infected compartments. Instead of analyzing the whole population, it isolates the infected groups and divides the rate of "new infections" by the rate of "recoveries" to easily calculate $R_0$.}. Their matrix approach standardized the computation of $R_0$, which we directly apply in our local stability analysis \hyperref[vandendriessche2002]{(van den Driessche and Watmough 30)}.

When analyzing the long-term, global behavior of these diseases, early researchers struggled to prove that populations would safely reach a stable equilibrium without crashing. C. Connell McCluskey solved this by proving that a specific mathematical construct—the Volterra-type Lyapunov function\footnote{\textbf{Volterra-type Lyapunov Function:} A mathematical "energy" scoring system used to prove a physical system will eventually settle down to a stable state. It cleverly utilizes natural logarithms to create an "infinite mathematical wall" that prevents population variables from ever dropping below zero.}—could perfectly map the non-linear spread of viruses \hyperref[mccluskey2010]{(McCluskey 2405)}. Furthermore, Hans J. Bremermann and Horst R. Thieme laid the groundwork for multi-strain dynamics by mathematically proving the "Competitive Exclusion Principle," which states that in a shared environment, the virus strain with the highest reproduction number will inevitably drive all competing strains to extinction \hyperref[bremermann1989]{(Bremermann and Thieme 180)}.

\begin{table}[H]
    \centering
    \resizebox{\textwidth}{!}{
    \begin{tabular}{|l|p{3cm}|p{3.5cm}|p{4cm}|p{4cm}|}
        \hline
        \textbf{Author(s)} & \textbf{Model Type} & \textbf{Local Stability Method} & \textbf{Global Stability Method} & \textbf{Key Contribution / Computational Results} \\
        \hline
        \hyperref[brauer2012]{Brauer \& Castillo-Chavez} & SIR / SVIR & Jacobian Linearization & Basic Phase Portraits & Established foundational models for vaccine leakage. \\
        \hline
        \hyperref[vandendriessche2002]{van den Driessche \& Watmough} & Complex Compartmental & Next-Generation Matrix & N/A & Standardized the calculation of $R_0$ for multi-group systems. \\
        \hline
        \hyperref[mccluskey2010]{McCluskey} & Non-linear SIR & N/A & Volterra-type Lyapunov & Proved global stability using logarithmic energy functions. \\
        \hline
        \hyperref[bremermann1989]{Bremermann \& Thieme} & Multi-Strain & N/A & Competitive Exclusion Theorem & Proved the mathematical extinction of weaker viral strains. \\
        \hline
        \textbf{Our Project} & \textbf{Multi-Strain SVIR} & \textbf{Next-Generation Matrix} & \textbf{Volterra-type Lyapunov \& Competitive Exclusion (Simulated)} & \textbf{Synthesizes stability theorems with 3-year computational tracking of Strain displacement.} \\
        \hline
    \end{tabular}
    }
    \caption{Summary of foundational literature compared to the current study.}
    \label{tab:lit_review}
\end{table}

\textbf{Contrast and Contribution:} 
While the literature provides a robust theoretical framework, these foundational texts primarily treat stability analysis and computational modeling as separate domains. For instance, McCluskey and Bremermann focus strictly on rigorous, abstract proofs of global stability without applied numerical simulations, while modern public health papers often rely heavily on computer simulations without analytically proving the underlying matrix stability. 
Our project bridges this gap. We synthesize van den Driessche and Watmough's NGM method to find the local outbreak thresholds and adapt McCluskey's Volterra structures to prove the global stability of the disease-free state. Crucially, we then step beyond the abstract proofs to perform a highly specific computational sensitivity analysis. By simulating a 1095-day competition between a baseline strain ($\beta_1 = 0.60, \sigma_1 = 0.10$) and a highly evasive strain ($\beta_2 = 1.16, \sigma_2 = 0.60$), we physically visualize Bremermann and Thieme's Competitive Exclusion Principle in real-time. We not only calculate the exact tipping point where the evasive strain becomes viable ($\beta_2 \approx 0.073$), but we also computationally chart the phase transition where it actively dominates the baseline strain ($\beta_2 > 0.585$)—a synthesis of pure mathematics and applied computational epidemiology that builds directly upon these classic works.


% ==========================================
% 4. LOCAL STABILITY ANALYSIS
% ==========================================
\section{Local Stability Analysis}

We define a system to be ``locally stable" if it can survive a tiny disturbance. Imagine a marble resting at the bottom of a bowl; if you flick it slightly, gravity pulls it back to the center. In epidemiology, that center point is a completely healthy population, and the ``flick" is the sudden introduction of a few infected tourists. Local stability analysis allows us to analytically prove whether the virus will die out (the marble returns to the center) or explode into a pandemic (the marble escapes the bowl).

\subsection{Disease-Free Equilibrium (DFE)}
The first step is to find the exact coordinates of that resting center point, known as the \textbf{Disease-Free Equilibrium (DFE)}. This represents a steady state where the population is completely free of both viral strains ($I_1 = 0$ and $I_2 = 0$).

From our susceptible equation (\ref{eq:S}), setting all time derivatives and infected variables to zero yields:
\begin{align*}
\Lambda - (\nu + \mu)S &= 0 \\
S^* &= \frac{\Lambda}{\nu + \mu}
\end{align*}
This confirms that in a healthy world, the susceptible population ($S^*$) is a simple balance between the birth rate ($\Lambda$) and the rate at which people leave the group through vaccination or natural mortality ($\nu + \mu$).

Substituting $S^*$ into the vaccinated equation (\ref{eq:V}):
\begin{align*}
\nu S^* - \mu V &= 0 \\
V^* &= \frac{\nu}{\mu} S^* = \frac{\nu \Lambda}{\mu(\nu + \mu)}
\end{align*}
Finally, with no active infections, the recovered compartment naturally empties out to zero ($R^* = 0$). Thus, the precise coordinates for our Disease-Free Equilibrium, denoted as $E_0$, are:
\begin{equation}
E_0 = \left( \frac{\Lambda}{\nu + \mu}, \frac{\nu \Lambda}{\mu(\nu + \mu)}, 0, 0, 0 \right)
\end{equation}

\subsection{Jacobian Matrix and Local Stability}
To understand why we evaluate the Jacobian at the Disease-Free Equilibrium ($E_0$), we can apply a first-order Taylor series expansion to linearize our non-linear system around that point. If we define a very small perturbation (e.g., a localized outbreak) as $\Delta x(t)$, its rate of change is approximately governed by the Jacobian matrix $J$ evaluated at the equilibrium:
\begin{equation}
\frac{d\Delta x}{dt} \approx J(E_0) \Delta x
\end{equation}
Because this is now a system of linear differential equations, its general solution takes the form of an exponential function driven by the matrix's eigenvalues ($\lambda$) and eigenvectors ($v$):
\begin{equation}
\Delta x(t) = c_1 e^{\lambda_1 t} v_1 + c_2 e^{\lambda_2 t} v_2 + \dots
\end{equation}
From a computational perspective, if the real part of every eigenvalue is strictly negative ($\text{Re}(\lambda_i) < 0$), the perturbation $\Delta x(t)$ exponentially decays to $0$ as $t \to \infty$. The system naturally terminates the outbreak and returns to $E_0$.

Because our system of equations is non-linear, we evaluate its stability by linearizing it around the $E_0$ equilibrium point. We construct the general Jacobian matrix ($J$) by taking the partial derivatives of our five differential equations with respect to our five variables ($S, V, I_1, I_2, R$):

\begin{equation}
\resizebox{\textwidth}{!}{
$J = \begin{bmatrix}
-(\beta_1 I_1 + \beta_2 I_2) - (\nu + \mu) & 0 & -\beta_1 S & -\beta_2 S & 0 \\
\nu & -(\sigma_1 \beta_1 I_1 + \sigma_2 \beta_2 I_2) - \mu & -\sigma_1 \beta_1 V & -\sigma_2 \beta_2 V & 0 \\
\beta_1 I_1 & \sigma_1 \beta_1 I_1 & (\beta_1 S + \sigma_1 \beta_1 V) - (\gamma_1 + \mu) & 0 & 0 \\
\beta_2 I_2 & \sigma_2 \beta_2 I_2 & 0 & (\beta_2 S + \sigma_2 \beta_2 V) - (\gamma_2 + \mu) & 0 \\
0 & 0 & \gamma_1 & \gamma_2 & -\mu
\end{bmatrix}$
}
\end{equation}

Evaluating this matrix at $E_0$ significantly simplifies the system, as $I_1$ and $I_2$ become zero. Because rows 3 and 4 contain zeros in the $S$, $V$, and $R$ columns, the variables conceptually decouple into a block-triangular format. We can visually demonstrate this structural decoupling by partitioning the matrix, isolating the infected block (center) from the non-infected compartments:

\begin{equation}
\resizebox{\textwidth}{!}{
$J(E_0) = \left[
\begin{array}{cc|cc|c}
-(\nu + \mu) & 0 & -\beta_1 S^* & -\beta_2 S^* & 0 \\
\nu & -\mu & -\sigma_1 \beta_1 V^* & -\sigma_2 \beta_2 V^* & 0 \\
\hline
0 & 0 & (\beta_1 S^* + \sigma_1 \beta_1 V^*) - (\gamma_1 + \mu) & 0 & 0 \\
0 & 0 & 0 & (\beta_2 S^* + \sigma_2 \beta_2 V^*) - (\gamma_2 + \mu) & 0 \\
\hline
0 & 0 & \gamma_1 & \gamma_2 & -\mu
\end{array}
\right]$
}
\end{equation}

Because the matrix is block-triangular, its spectrum is simply the union of the eigenvalues of the decoupled blocks. To determine the eigenvalues corresponding to the non-infected compartments, we remove the infected rows and columns (row 3, column 3 and row 4, column 4). This yields the reduced $3 \times 3$ sub-matrix for $S, V,$ and $R$:

\begin{equation}
J_{SVR}(E_0) = \begin{bmatrix}
-(\nu + \mu) & 0 & 0 \\
\nu & -\mu & 0 \\
0 & 0 & -\mu
\end{bmatrix}
\end{equation}

Because $J_{SVR}(E_0)$ is a lower-triangular matrix, its eigenvalues are found directly on its main diagonal:
\begin{equation}
\lambda_1 = -(\nu + \mu), \quad \lambda_2 = -\mu, \quad \lambda_5 = -\mu
\end{equation}
Since the real-world parameters for vaccination ($\nu$) and natural mortality ($\mu$) are strictly positive, $\lambda_1, \lambda_2,$ and $\lambda_5$ are definitively negative. 

Therefore, the stability of the entire system depends exclusively on the remaining two eigenvalues extracted from the decoupled $2 \times 2$ infected block:
\begin{align}
\lambda_3 &= (\beta_1 S^* + \sigma_1 \beta_1 V^*) - (\gamma_1 + \mu) \label{eq:lam3} \\
\lambda_4 &= (\beta_2 S^* + \sigma_2 \beta_2 V^*) - (\gamma_2 + \mu) \label{eq:lam4}
\end{align}
For $E_0$ to be locally asymptotically stable, we require both $\lambda_3 < 0$ and $\lambda_4 < 0$. These threshold conditions lead directly to the calculation of the Basic Reproduction Number.

\subsection{Basic Reproduction Number ($R_0$)}
To formally extract $R_0$, we apply the Next-Generation Matrix (NGM) method \hyperref[vandendriessche2002]{(van den Driessche and Watmough 30)}. The NGM method isolates the infected compartments and mathematically decomposes their dynamics into two distinct vectors: the rate of entirely new infections ($\mathcal{F}$) and the rate of transitions/recoveries out of the infected state ($\mathcal{V}$).

\begin{equation}
\mathcal{F} = \begin{bmatrix} (\beta_1 S + \sigma_1 \beta_1 V)I_1 \\ (\beta_2 S + \sigma_2 \beta_2 V)I_2 \end{bmatrix}, \quad
\mathcal{V} = \begin{bmatrix} (\gamma_1 + \mu)I_1 \\ (\gamma_2 + \mu)I_2 \end{bmatrix}
\end{equation}

By taking the Jacobian of $\mathcal{F}$ and $\mathcal{V}$ with respect to $I_1$ and $I_2$, and evaluating them at the DFE, we generate the corresponding matrices $F$ and $V$:

\begin{equation}
F = \begin{bmatrix}
\beta_1 S^* + \sigma_1 \beta_1 V^* & 0 \\
0 & \beta_2 S^* + \sigma_2 \beta_2 V^*
\end{bmatrix}, \quad
V = \begin{bmatrix}
\gamma_1 + \mu & 0 \\
0 & \gamma_2 + \mu
\end{bmatrix}
\end{equation}

The Next-Generation Matrix, which calculates the ratio of new infections generated per single recovery, is defined as $K = F V^{-1}$:

\begin{equation}
K = \begin{bmatrix}
\frac{\beta_1 S^* + \sigma_1 \beta_1 V^*}{\gamma_1 + \mu} & 0 \\
0 & \frac{\beta_2 S^* + \sigma_2 \beta_2 V^*}{\gamma_2 + \mu}
\end{bmatrix}
\end{equation}

While calculating $K = F V^{-1}$ appears abstract, it is a highly logical matrix operation when broken down into discrete steps. 
\begin{itemize}
    \item \textbf{The $F$ Matrix} represents the rate at which current infected states generate \textit{new} infections.
    \item \textbf{The $V$ Matrix} represents the rate at which individuals leave the infected state. Therefore, its inverse, $V^{-1}$, represents the \textbf{expected time} an individual remains infectious.
\end{itemize}
When we multiply them, we are simply calculating:
\begin{equation*}
(\text{Rate of new infections}) \times (\text{Average time spent infected}) = \text{Total Secondary Infections}
\end{equation*}
Thus, the Next-Generation Matrix $K$ directly computes the expected number of new infections generated by a single case.

The basic reproduction number $R_0$ is mathematically defined as the dominant eigenvalue (spectral radius) of $K$. Because $K$ is a diagonal matrix, its eigenvalues are simply its diagonal entries. Thus, we successfully extract the strain-specific reproduction numbers:

\begin{align}
R_{0,1} &= \frac{\beta_1 S^* + \sigma_1 \beta_1 V^*}{\gamma_1 + \mu} \\
R_{0,2} &= \frac{\beta_2 S^* + \sigma_2 \beta_2 V^*}{\gamma_2 + \mu}
\end{align}

The reproduction number for the entire system is $R_0 = \max(R_{0,1}, R_{0,2})$. If $R_0 < 1$, the Disease-Free Equilibrium is locally asymptotically stable, guaranteeing that any initial outbreak will safely diminish.


% ==========================================
% 5. GLOBAL STABILITY ANALYSIS
% ==========================================
\section{Global Stability Analysis}

While Local Stability Analysis proves that a healthy population can survive a \textit{small} disturbance (like a few infected tourists), it does not guarantee safety against a massive shock. Global Stability Analysis tests whether the system will eventually return to the Disease-Free Equilibrium (DFE) regardless of the starting conditions—even if 90\% of the population is suddenly infected. To return to our previous analogy: if local stability is a marble returning to the center after a tiny flick, global stability proves that no matter where in the room you drop the marble, the floor is slanted to force it back into the bowl.

\subsection{The Volterra-Type Lyapunov Function}
Because we cannot perfectly solve our non-linear differential equations to see exactly where the population will be in the distant future, we use a mathematical ``cheat code" known as a \textbf{Lyapunov Function} ($L$). 

Think of $L$ as an artificial ``energy" or ``error" score that measures exactly how far our current population is from the perfect, healthy equilibrium ($E_0$). If we can prove that this energy score is strictly decreasing over time ($\frac{dL}{dt} \le 0$), the laws of physics and calculus dictate that the system must eventually run out of energy and come to a complete stop exactly at $E_0$.

To construct this for our model, we use a Volterra-type Lyapunov function, heavily utilizing natural logarithms. We define our function as the sum of the energies of our four active compartments:
\begin{equation} \label{eq:lyapunov}
L = \left( S - S^* - S^* \ln \frac{S}{S^*} \right) + \left( V - V^* - V^* \ln \frac{V}{V^*} \right) + I_1 + I_2
\end{equation}

\textbf{Core Assumption:} We assume that the population compartments ($S, V$) are strictly positive ($> 0$). This is biologically logical, but it is also mathematically enforced by the shape of the logarithmic terms in Equation \ref{eq:lyapunov}. As $S$ or $V$ drops closer to zero, the natural logarithm approaches negative infinity. Because we are subtracting that logarithm, the overall "energy" score $L$ shoots up to positive infinity. This creates an invisible, infinitely high mathematical wall that physically repels the population from ever dropping below zero.

\subsection{Derivative and Cancellation}
To prove the system loses energy over time, we must take the time derivative of the entire Lyapunov function ($\frac{dL}{dt}$) and prove it is negative. By applying the chain rule to the logarithmic terms, the derivative separates into:
\begin{equation}
\frac{dL}{dt} = \left(1 - \frac{S^*}{S}\right)\frac{dS}{dt} + \left(1 - \frac{V^*}{V}\right)\frac{dV}{dt} + \frac{dI_1}{dt} + \frac{dI_2}{dt}
\end{equation}

Next, we substitute our original ODEs (Equations \ref{eq:S} - \ref{eq:I2}) into these derivatives. At first glance, this creates a massive algebraic mess. However, the brilliance of the Volterra logarithmic structure is that it perfectly cancels out the non-linear disease transmission terms. Let us isolate and observe the terms associated strictly with Strain 1 ($I_1$):

From the susceptible ($S$) expansion:
\begin{equation*}
\left(1 - \frac{S^*}{S}\right)(-\beta_1 S I_1) = -\beta_1 S I_1 + \beta_1 S^* I_1
\end{equation*}
From the vaccinated ($V$) expansion:
\begin{equation*}
\left(1 - \frac{V^*}{V}\right)(-\sigma_1 \beta_1 V I_1) = -\sigma_1 \beta_1 V I_1 + \sigma_1 \beta_1 V^* I_1
\end{equation*}
From the newly infected ($I_1$) expansion:
\begin{equation*}
\frac{dI_1}{dt} = \beta_1 S I_1 + \sigma_1 \beta_1 V I_1 - (\gamma_1 + \mu)I_1
\end{equation*}

When we add these three lines together, the active infection terms ($-\beta_1 S I_1$ and $+\beta_1 S I_1$) perfectly cancel each other out. The same happens for the vaccine transmission terms. We are left only with the equilibrium terms:
\begin{equation*}
\text{Sum of } I_1 \text{ terms} = \left[ \beta_1 S^* + \sigma_1 \beta_1 V^* - (\gamma_1 + \mu) \right] I_1
\end{equation*}

By factoring out the recovery/death rate $(\gamma_1 + \mu)$, we expose a deeply familiar ratio:
\begin{equation*}
(\gamma_1 + \mu) \left[ \frac{\beta_1 S^* + \sigma_1 \beta_1 V^*}{\gamma_1 + \mu} - 1 \right] I_1
\end{equation*}
The fraction inside the bracket is the exact definition of our Basic Reproduction Number ($R_{0,1}$). By repeating this exact identical cancellation process for the $S$, $V$, and $I_2$ terms, the massive derivative collapses into an incredibly elegant final form:

\begin{equation} \label{eq:lyapunov_final}
\frac{dL}{dt} = -\frac{(\nu + \mu)}{S}(S - S^*)^2 - \frac{\mu}{V}(V - V^*)^2 + (\gamma_1 + \mu)(R_{0,1} - 1)I_1 + (\gamma_2 + \mu)(R_{0,2} - 1)I_2
\end{equation}

\subsection{Global Stability Conclusion}
Observe the final equation (\ref{eq:lyapunov_final}). The first two terms, describing $S$ and $V$, are squared $(S - S^*)^2$. Since any real number squared is positive, and they are multiplied by negative fractions, these first two terms are \textbf{permanently negative}. 

The entire fate of the system therefore rests on the $I_1$ and $I_2$ terms. Because physical populations ($I_1, I_2$) and vital rates ($\gamma, \mu$) are strictly positive, the only way for these final terms to be negative is if $(R_{0,1} - 1) \le 0$ and $(R_{0,2} - 1) \le 0$. 

Thus, if the overall Basic Reproduction Number $R_0 \le 1$, then $\frac{dL}{dt} \le 0$ for all strictly positive population sizes.

From an algorithmic perspective, knowing that the system's ``error" or ``energy" never increases ($\le 0$) is good, but it raises a critical edge-case: what if the system hits a state where $\frac{dL}{dt} = 0$ but it is not at the exact target equilibrium? Could the system get stuck in an infinite loop or a local plateau? 

To solve this, we invoke \textbf{LaSalle's Invariance Principle}. This topological theorem acts as a strict termination proof. It states that an autonomous system will eventually settle into the largest ``invariant set" where $\frac{dL}{dt} = 0$. By examining our Lyapunov derivative, the only possible mathematical state where the derivative remains exactly zero forever is when there are no infected individuals ($I_1 = 0, I_2 = 0$) and the susceptible pools are at their exact baselines ($S = S^*, V = V^*$).

Therefore, by LaSalle's Invariance Principle, we conclude that the Disease-Free Equilibrium ($E_0$) is \textbf{Globally Asymptotically Stable}. The viruses will inevitably be eradicated, regardless of the initial outbreak size, because the system physically cannot halt anywhere else.

\subsection{Global Stability of the Endemic Equilibrium and Competitive Exclusion}
The rigorous proof above hinges on the assumption that $R_0 \le 1$. However, what happens if a highly infectious variant pushes $R_0 > 1$? 

Looking back at Equation \ref{eq:lyapunov_final}, if $R_0 > 1$, the infected terms become massive positive numbers, overpowering the negative $S$ and $V$ terms. Consequently, $\frac{dL}{dt} > 0$, meaning the ``energy" or ``distance" from the healthy state actively increases. The system physically runs away from the DFE because the DFE has become utterly unstable. The population must instead settle into a new, stable "bowl" known as the \textbf{Endemic Equilibrium (EE)}, where the virus circulates perpetually.

While this section explicitly derives the global stability of the DFE, establishing the Global Asymptotic Stability of the non-zero Endemic Equilibrium requires the construction of a substantially more complex, multi-dimensional Volterra Lyapunov functional centered directly on the endemic coordinates. The foundational work by C. Connell McCluskey established that such logarithmic functionals successfully prove the global stability of endemic states in non-linear compartmental models \hyperref[mccluskey2010]{(McCluskey 2405)}. 

Furthermore, regarding the multi-strain dynamics observed in our system, Hans J. Bremermann and Horst R. Thieme established the rigorous mathematical proof for the ``Competitive Exclusion Principle." Their theorem mathematically guarantees that in a strictly competitive multi-strain environment, the single strain maximizing the basic reproduction number will globally exclude all other strains, achieving a globally stable dominance equilibrium while driving its competitors to absolute zero \hyperref[bremermann1989]{(Bremermann and Thieme 180)}. 

Therefore, the multi-strain endemic stability and phase transitions we explore computationally in the subsequent sections of this report are fully supported by established global stability theorems in the epidemiological literature.


% ==========================================
% 6. COMPUTATIONAL SIMULATION
% ==========================================
\section{Computational Simulation}

While the global stability theorems in the previous section mathematically guarantee the final resting state of the population, they do not show us the turbulent journey required to get there. To bridge the gap between abstract algebra and real-world epidemiology, we utilize numerical integration (computational simulation) to track the exact daily changes of all five population compartments over a 180-day acute outbreak period. 

By applying parameters that mirror the biological realities of the COVID-19 pandemic—specifically the competition between a baseline strain (representing Delta) and a highly transmissible, vaccine-evasive mutant (representing Omicron)—we can visualize the Competitive Exclusion Principle in real time.

\subsection{Scenario A: Baseline Competition}
In our first simulation, we introduce both viral strains into a highly susceptible population with a standard, baseline vaccination rate. Strain 1 is highly infectious, but Strain 2 features a higher fundamental transmission rate ($\beta_2 > \beta_1$) and a severe vaccine leakage factor ($\sigma_2 > \sigma_1$), meaning it can easily infect vaccinated individuals. 

\begin{figure}[H]
    \centering
    \includegraphics[width=\textwidth]{Scenario-A.jpeg}
    \caption{Time-series simulation over 180 days demonstrating the competitive exclusion of Strain 1 (red) by Strain 2 (purple).}
    \label{fig:scenario_a}
\end{figure}

As shown in Figure \ref{fig:scenario_a}, Strain 1 (red dashed line) initially attempts to spark an outbreak. However, Strain 2 (purple dashed line) explodes through the population, reaching a peak prevalence of over 70\% by day 10. 

Because both viruses rely on the exact same susceptible pool (blue line) for "fuel," Strain 2's rapid expansion aggressively starves Strain 1. Before Strain 1 can establish a foothold, the susceptible population is entirely depleted, and the recovered population (green line) surges to nearly 100\%. This serves as a perfect visual proof of Bremermann and Thieme's Competitive Exclusion Principle: the strain with the higher reproductive advantage does not just coexist; it actively drives the lesser strain into early extinction.

% ==========================================
% 7. SENSITIVITY ANALYSIS
% ==========================================
\section{Sensitivity Analysis}

Having established the baseline dynamics, we now alter specific parameters to test the limits of the system. This allows us to answer critical public health questions: Can a massive vaccination campaign stop the highly evasive mutant? And exactly how infectious does a new strain need to be to take over the population?

\subsection{Scenario B: Vaccination Surge}
In Scenario B, we simulate a massive public health intervention by drastically increasing the vaccination rate ($\nu$). 

\begin{figure}[H]
    \centering
    \includegraphics[width=\textwidth]{Scenario-B.jpeg}
    \caption{Outbreak dynamics under a sharply increased vaccination rate.}
    \label{fig:scenario_b}
\end{figure}

Observing Figure \ref{fig:scenario_b}, the intervention initially appears successful. The vaccinated compartment (orange line) spikes to over 20\% in the first few days, a significant increase compared to Scenario A. Consequently, the susceptible pool (blue line) crashes even faster as people are moved to relative safety.

However, the peak of Strain 2 (purple line) remains remarkably high, cresting near 70\%. Why did the vaccination surge fail to suppress the outbreak? The answer lies in the vaccine leakage parameter ($\sigma_2$). Because Strain 2 is mathematically defined to evade vaccine-induced immunity, the newly vaccinated individuals simply become an alternative fuel source for the virus. This mathematically proves a grim epidemiological reality: against a highly evasive strain, increasing the rate of a "leaky" vaccine accelerates the suppression of the weaker strain but does little to blunt the peak of the evasive dominant strain.

\subsection{Scenario C: Sensitivity of Strain 2 to Transmission Rate}
Finally, we want to map the exact threshold at which Strain 2 achieves this absolute dominance. Rather than graphing time on the x-axis, Figure \ref{fig:scenario_c} plots the varying transmission rate of Strain 2 ($\beta_2$) against the maximum peak size the outbreak achieves. 

\begin{figure}[H]
    \centering
    \includegraphics[width=\textwidth]{Scenario-C.jpeg}
    \caption{Logistic phase transition showing the peak prevalence of Strain 2 relative to its transmission rate ($\beta_2$).}
    \label{fig:scenario_c}
\end{figure}

This sensitivity curve illustrates a dramatic "phase transition" in the virus's viability. 
\begin{itemize}
    \item \textbf{Extinction Phase ($\beta_2 < 0.3$):} If the transmission rate is too low, the peak prevalence is practically zero. The virus is competitively excluded by Strain 1 or naturally dies out.
    \item \textbf{The Tipping Point ($0.4 < \beta_2 < 0.8$):} As the transmission rate crosses the critical reproduction threshold, we see an exponential explosion in peak prevalence. This steep slope represents the exact mathematical window where Strain 2 overtakes Strain 1.
    \item \textbf{Saturation Phase ($\beta_2 > 1.0$):} At extreme transmission rates, the curve flattens out, asymptoting near 80\%. At this stage, the virus is moving so fast that it infects nearly everyone available before they can naturally recover or be vaccinated. 
\end{itemize}

Ultimately, this sensitivity analysis proves that competitive dominance is not a gradual shift; it is a rapid, non-linear threshold event strictly dictated by the mathematical parameters of transmission and evasion.


% ==========================================
% 8. CONCLUSION AND FUTURE WORK
% ==========================================
\section{Conclusion and Future Work}

\subsection{Conclusion}
This project successfully constructed and analyzed a multi-strain SVIR epidemic model to explore the competitive dynamics between a baseline viral strain and a highly transmissible, vaccine-evasive mutation. By bridging rigorous mathematical proofs with dynamic computational simulations, we were able to map both the theoretical limitations and the real-world trajectories of the viruses.

Analytically, we determined the exact threshold for an outbreak by calculating the Basic Reproduction Number ($R_0$) using the Next-Generation Matrix method. We then proved the Global Asymptotic Stability of the Disease-Free Equilibrium by constructing a Volterra-type Lyapunov function, mathematically guaranteeing that if $R_0 \le 1$, the virus will eventually be eradicated regardless of the initial outbreak size. 

Computationally, we pushed the model beyond this safe equilibrium to simulate a real-world pandemic scenario ($R_0 > 1$) over a 180-day period. Our numerical integration vividly demonstrated the Competitive Exclusion Principle. We proved that when two viruses compete for the same susceptible population, they do not peacefully coexist. Instead, the highly evasive strain acts as a biological vacuum, rapidly consuming the available susceptible and vaccinated hosts, actively starving the baseline strain, and driving it into early extinction. 

Ultimately, our sensitivity analysis revealed a sobering public health reality: while mass vaccination campaigns are highly effective at suppressing a baseline strain, they cannot stop a highly evasive mutation from achieving peak endemicity. The dominance of a new variant is not a gradual shift, but a brutal, non-linear phase transition dictated strictly by its transmission and evasion parameters.

\subsection{Future Work}
While our model successfully captures the mechanics of variant displacement, mathematical epidemiology is an ever-evolving field. To further align this model with the complex realities of global pandemics, future research could expand upon our framework in several critical directions:

\begin{itemize}
    \item \textbf{Waning Immunity:} Our current model assumes that individuals in the Recovered ($R$) and Vaccinated ($V$) compartments maintain their immunological status indefinitely. Future iterations should introduce a decay rate, allowing recovered and vaccinated individuals to flow back into the Susceptible ($S$) compartment over time, simulating the need for booster shots.
    \item \textbf{Behavioral and Policy Dynamics:} The transmission rates ($\beta_1, \beta_2$) in our system remain constant. Incorporating time-dependent variables to simulate government interventions—such as lockdowns, mask mandates, or social distancing—would allow us to study how human behavior actively alters the competitive exclusion threshold.
    \item \textbf{Age-Structured Compartments:} Vaccines do not uniformly protect all age groups, and natural mortality rates ($\mu$) differ vastly between youth and elderly populations. Sub-dividing the compartments into a matrix of age groups would yield highly targeted, demographic-specific stability analysis.
    \item \textbf{Spatial and Network Modeling:} We assumed a perfectly mixed, homogenous population. Expanding this to a network-based model where nodes represent different cities or countries would allow researchers to track how international travel and localized vaccination rates affect the global spread of competing strains.
\end{itemize}

% ==========================================
% WORKS CITED (MLA Format)
% ==========================================
\newpage
\renewcommand{\refname}{Works Cited}
\begin{thebibliography}{9}

\bibitem{brauer2012} \label{brauer2012}
Brauer, Fred, and Carlos Castillo-Chavez. \textit{Mathematical Models in Population Biology and Epidemiology}. 2nd ed., Springer, 2012.

\bibitem{bremermann1989} \label{bremermann1989}
Bremermann, Hans J., and Horst R. Thieme. ``A Competitive Exclusion Principle for Pathogen Strains.'' \textit{Journal of Mathematical Biology}, vol. 27, no. 2, 1989, pp. 179-190.

\bibitem{mccluskey2010} \label{mccluskey2010}
McCluskey, C. Connell. ``Complete Global Stability for an SIR Epidemic Model with Delay---Distributed or Discrete.'' \textit{Nonlinearity}, vol. 23, no. 9, 2010, p. 2405.

\bibitem{vandendriessche2002} \label{vandendriessche2002}
van den Driessche, Pauline, and James Watmough. ``Reproduction Numbers and Sub-Threshold Endemic Equilibria for Compartmental Models of Disease Transmission.'' \textit{Mathematical Biosciences}, vol. 180, no. 1-2, 2002, pp. 29-48.

\end{thebibliography}

\end{document}
